\section{Giới thiệu}

\subsection{Lý do chọn đề tài}

Việc trồng cây trong nhà ngày càng phổ biến, đặc biệt tại các khu vực đô thị có
không gian hạn chế. Tuy nhiên, người chăm sóc cây không phải lúc nào cũng có thể
theo dõi độ ẩm đất, nhiệt độ hoặc tưới cây đúng thời điểm. Lịch làm việc, những
chuyến đi xa và sự thay đổi của thời tiết đều có thể làm gián đoạn quá trình chăm
sóc, ảnh hưởng đến sức khỏe của cây.

Các hệ thống nông nghiệp quy mô lớn thường có chi phí và yêu cầu lắp đặt vượt quá
nhu cầu của một hộ gia đình. Từ thực tế đó, nhóm lựa chọn xây dựng SmartGarden,
một mô hình IoT nhỏ gọn có khả năng thu thập dữ liệu môi trường, điều khiển thiết
bị và hỗ trợ chăm sóc cây từ xa. Đề tài đồng thời tạo điều kiện để nhóm kết hợp
kiến thức về cảm biến, mạch điều khiển, truyền thông mạng, dịch vụ đám mây và phát
triển phần mềm trong cùng một sản phẩm.

\subsection{Mục tiêu đồ án}

Mục tiêu của đồ án là xây dựng một mô hình SmartGarden hoàn chỉnh từ tầng thiết
bị đến tầng ứng dụng. Hệ thống sử dụng ESP32 để đọc cảm biến và điều khiển tải;
back-end tiếp nhận dữ liệu, quản lý nghiệp vụ và lưu trữ; Website cùng Telegram
Bot cung cấp giao diện theo dõi và điều khiển cho người dùng.

Các mục tiêu cụ thể gồm:

\begin{itemize}
    \item \textbf{Phần cứng:} lắp ráp bộ điều khiển ESP32, các cảm biến môi
    trường, relay, máy bơm, quạt và đèn; bảo đảm relay ở trạng thái an toàn khi
    khởi động và giới hạn thời gian chạy của máy bơm.
    \item \textbf{Kết nối:} trao đổi dữ liệu giữa ESP32 và back-end qua MQTT/TLS,
    đồng thời duy trì cơ chế kết nối lại khi Wi-Fi hoặc máy chủ MQTT bị gián đoạn.
    \item \textbf{Phần mềm:} xây dựng REST API, Website và Telegram Bot dùng
    chung lớp nghiệp vụ, hỗ trợ xác thực, phân quyền, giám sát, điều khiển, lập
    lịch, cảnh báo và phân tích dữ liệu.
    \item \textbf{Lưu trữ:} sử dụng Firebase Authentication và Realtime Database
    để quản lý tài khoản, dữ liệu cảm biến, trạng thái thiết bị và lịch sử vận
    hành.
\end{itemize}

Các tiêu chí vận hành được xác định cho phiên bản hiện tại như sau:

\begin{itemize}
    \item ESP32 đọc và gửi dữ liệu cảm biến theo chu kỳ khoảng 1 giây; trạng thái
    thiết bị và tín hiệu duy trì kết nối được gửi khoảng 5 giây một lần.
    \item Người dùng có thể điều khiển máy bơm, quạt và đèn từ Website hoặc
    Telegram sau khi được xác thực.
    \item Hệ thống hỗ trợ ba chế độ vận hành: thủ công, theo ngưỡng và theo lịch.
    Máy bơm được giới hạn thời gian chạy liên tục ở mức 180 giây.
    \item Các tình huống đất khô, nhiệt độ cao, máy bơm chạy quá lâu và thiết bị
    mất kết nối được phát hiện để tạo cảnh báo.
    \item Việc truy cập từ Internet được thực hiện khi Website và back-end được
    triển khai trên máy chủ; trong giai đoạn phát triển, hệ thống có thể chạy
    trong mạng nội bộ.
    \item Chi phí phần cứng được kiểm soát theo danh sách linh kiện trình bày ở
    Mục 4; chi phí dịch vụ đám mây không nằm trong phạm vi thống kê này.
\end{itemize}

\subsection{Giới thiệu sản phẩm SmartGarden}

SmartGarden là mô hình quản lý vườn thông minh dành cho nhu cầu chăm sóc cây tại
nhà. Các cảm biến cung cấp dữ liệu nhiệt độ, độ ẩm không khí, độ ẩm đất và ánh
sáng. ESP32 xử lý dữ liệu tại thiết bị, điều khiển máy bơm, quạt và đèn, đồng
thời trao đổi với back-end qua MQTT. Back-end đảm nhiệm xác thực, lưu trữ, cảnh
báo, quản lý cấu hình và cung cấp REST API cho Website và Telegram Bot.

Hệ thống được tổ chức thành ba lớp:

\begin{itemize}
    \item \textbf{Lớp thiết bị:} ESP32, cảm biến môi trường, relay và các thiết
    bị chấp hành.
    \item \textbf{Lớp dịch vụ:} Flask API, Firebase, máy chủ MQTT, bộ xử lý lịch,
    dịch vụ cảnh báo và dịch vụ phân tích AI.
    \item \textbf{Lớp ứng dụng:} Website phục vụ giám sát và quản trị, cùng
    Telegram Bot phục vụ thao tác nhanh và nhận thông báo.
\end{itemize}

Trong phạm vi đồ án, thành phẩm được triển khai cho một chậu cây với một máy bơm
và một nguồn nước nhỏ. Quy mô này đủ để minh họa luồng thu thập dữ liệu, điều
khiển, tự động hóa và cảnh báo của hệ thống. Kiến trúc phần mềm được tách thành
các mô-đun để thuận tiện mở rộng thêm thiết bị hoặc khu vực trồng trong các phiên
bản tiếp theo.
